\chapter{Plans for Thesis B and C}
\label{ch:projectplan}

This chapter outlines the planned work for Thesis~B and Thesis~C over an
eighteen--week period.\footnote{The exact calendar weeks may shift slightly
with UNSW term dates, but the relative ordering of tasks will be preserved.}
The project is organised around a single coherent case study: a
containerised microservice whose behaviour under varying load and CPU limits
is measured, modelled and used to teach capacity planning.  The work
progresses from consolidating the existing experimental platform, through
single-- and multi--server discrete--event simulation (DES), to hybrid
ML--DES modelling, evaluation, and the development of teaching artefacts for
COMP9334.

\section{Work Breakdown}

The work for Thesis~B and~C is grouped into the following major strands:

\begin{enumerate}
  \item \textbf{Stabilise and extend the experimental platform.}
  Finalise the Docker--based microservice, Vegeta load--generation scripts,
  and the cAdvisor/Prometheus/Grafana monitoring stack.  Automate
  experiments to collect traces of arrival rates, service times, CPU
  utilisation and tail latency for a representative range of loads and CPU
  quotas.

  \item \textbf{Single--server DES modelling and calibration.}
  Build a baseline DES model (single core, single threaded service) that
  reproduces measured behaviour under fixed load.  Calibrate the model
  using empirical $\lambda$ and service--time samples, and validate it
  against observed mean response time and utilisation.

  \item \textbf{Multi--server and multicore DES extensions.}
  Extend the simulator to multi--server (M/M/$c$ or M/G/$c$ style) cases,
  modelling multiple service threads or container replicas.  Incorporate
  simple models of CPU throttling to capture noisy--neighbour effects at
  higher utilisation.

  \item \textbf{Machine--learning models for service time and workload.}
  Train ML models to (i) learn empirical service--time distributions as a
  function of CPU quota and request rate (e.g.\ KDE, Gaussian mixture
  models, gradient--boosted regressors), and (ii) forecast arrival rates
  $\lambda(t)$ from historical monitoring data (e.g.\ ARIMA, gradient
  boosting, or lightweight sequence models).  These models provide
  distributional parameters and workload trajectories to feed into the DES
  engine.

  \item \textbf{Hybrid ML--DES integration and evaluation.}
  Integrate the ML components with the DES simulator so that the simulator
  operates on ML--estimated service--time distributions and forecasted
  arrivals.  Evaluate how well the hybrid model predicts mean latency,
  utilisation and tail behaviour (p95/p99) under different load and resource
  scenarios, and compare against purely analytical and purely ML baselines.

  \item \textbf{Educational artefacts and thesis consolidation.}
  Design teaching materials (lab scripts, small exercises and
  visualisations) that allow COMP9334 students to reproduce a subset of the
  experiments on their own machines.  In parallel, draft and refine the
  remaining thesis chapters (results, discussion and conclusions), ensuring
  that the experimental and modelling work is clearly tied back to the
  pedagogical goals.
\end{enumerate}

\section{Eighteen--Week Gantt Chart}

The high--level schedule is shown in
Figure~\ref{fig:gantt-thesis-bc}.  The timeline is divided into eighteen
weeks, covering both Thesis~B and Thesis~C.  Some activities deliberately
overlap: for example, data collection continues while the initial DES model
is being developed, and writing begins before all experimental work is
complete.

\begin{figure}[h]
  \centering
  \begin{ganttchart}[
      hgrid,
      vgrid,
      x unit=0.45cm,
      y unit chart=0.6cm,
      bar height=0.5,
      bar label font=\small,
      milestone label font=\small\bfseries
    ]{1}{18}
    \gantttitle{Weeks}{18} \\
    \gantttitlelist{1,...,18}{1} \\

    % Platform + data collection
    \ganttgroup{Platform and data}{1}{3} \\
    \ganttbar{Platform consolidation}{1}{2} \\
    \ganttbar{Baseline data collection}{2}{3} \\

    % DES work
    \ganttgroup{DES modelling}{3}{7} \\
    \ganttbar{Single--server DES model}{3}{5} \\
    \ganttbar{Multi--server / multicore DES}{5}{7} \\

    % ML models
    \ganttgroup{ML models}{7}{11} \\
    \ganttbar{Service--time modelling}{7}{10} \\
    \ganttbar{Workload forecasting}{9}{11} \\

    % Hybrid + eval
    \ganttgroup{Hybrid ML--DES \& evaluation}{11}{16} \\
    \ganttbar{ML--DES integration}{11}{14} \\
    \ganttbar{Evaluation \& analysis}{14}{16} \\

    % Teaching + write up
    \ganttgroup{Teaching artefacts \& write--up}{16}{18} \\
    \ganttbar{COMP9334 lab materials}{16}{17} \\
    \ganttbar{Final thesis consolidation}{16}{18} \\

    \ganttmilestone{Hybrid model stable}{14}
  \end{ganttchart}
  \caption{Planned timeline for Thesis~B and Thesis~C (18 weeks).}
  \label{fig:gantt-thesis-bc}
\end{figure}


This schedule leaves limited but intentional slack near the end of the
project (Weeks~16--18) to absorb unexpected implementation issues, repeat
key experiments, and polish both the thesis document and the accompanying
teaching materials.
